\documentclass[10pt,a4paper]{article}
\usepackage{amsmath,amssymb,amsthm,graphicx}
\usepackage{amsmath,amssymb,amsthm,hyperref,xfrac}
\providecommand{\eg}{e.g.\ }
\providecommand{\ie}{i.e.\ }
\usepackage[scr=rsfs,cal=euler]{mathalfa}
\multlinegap0pt
\hyphenation{accu-rate applied assump-tion assump-tions decom-pos-ing defined esti-mate esti-mates exten-ded hypothe-ses para-me-ter para-me-ters remain-der required}
\newenvironment{enumeratei}
{\bgroup\def\theenumi{\roman{enumi}}\def\theenumii{\arabic{enumii}}\begin{enumerate}}
{\end{enumerate}\egroup}
\newenvironment{enumeratea}
{\bgroup\def\theenumi{\alph{enumi}}\begin{enumerate}}
{\end{enumerate}\egroup}

\newcommand{\Psfrac}[2]{(\sfrac{#1}{#2})}
\newcommand{\psfrac}[2]{\sfrac{(#1)}{#2}}
\newcommand{\spfrac}[2]{\sfrac{#1}{(#2)}}
\newcommand{\pspfrac}[2]{\sfrac{(#1)}{(#2)}}
\newcommand{\shalf}{\sfrac{1}{2}}

\newcommand{\bnorm}[1]{\bigl\|#1\bigr\|}
\newcommand{\Bnorm}[1]{\Bigl\|#1\Bigr\|}
\newcommand{\babs}[1]{\bigl|#1\bigr|}
\newcommand{\Babs}[1]{\Bigl|#1\Bigr|}

\AtBeginDocument{%
\DeclareMathOperator{\sinn}{sin}\let\sin\sinn
\DeclareMathOperator{\coos}{cos}\let\cos\coos
}

%\usepackage{enumitem}
\usepackage{physics}
%\usepackage[foot]{amsaddr}
\usepackage{bm}
\usepackage{tikz}
\usepackage{mathtools}
\usepackage{scalerel,stackengine}
\stackMath
\usetikzlibrary{arrows.meta}
\usepackage{subcaption}
%\usepackage[style=numeric]{biblatex}
%\addbibresource{main.bib}
%\setlength{\bibitemsep}{0.3\baselineskip}

\newenvironment{mycases}
{\begin{dcases*}}
{\end{dcases*}}
\theoremstyle{plain}
\newtheorem{theorem}{Theorem}[section]
\newtheorem{lemma}[theorem]{Lemma}
\newtheorem{definition}[theorem]{Definition}
\newtheorem{proposition}[theorem]{Proposition}
\newtheorem{corollary}[theorem]{Corollary}
\newtheorem*{corollary*}{Corollary}
\theoremstyle{remark}
\newtheorem{remark}[theorem]{Remark}
\counterwithout{figure}{section}
\counterwithout{equation}{section}

\newcommand{\br}[1]{\left(#1\right)}
\newcommand{\bbr}[1]{\bigl(#1 \bigr)}
\newcommand{\Bbr}[1]{\Bigl(#1 \Bigr)}
\newcommand{\bbbr}[1]{\biggl(#1 \biggr)}
\newcommand{\tpd}{\frac{1}{(2\pi)^{d}}}
\newcommand\pig[1]{\scalerel*[2.35pt]{\big#1}{%
\ensurestackMath{\addstackgap[1.2pt]{\big#1}}}}
\newcommand\pigl[1]{\mathopen{\pig{#1}}}
\newcommand\pigr[1]{\mathclose{\pig{#1}}}
\newcommand{\ak}{\abs{k}}
\newcommand{\ap}{\abs{p}}
\newcommand{\aal}{\abs{\alpha}}
\newcommand{\sbbr}[1]{\bigl[#1\bigr]}
\newcommand{\sBbr}[1]{\Bigl[#1\Bigr]}
\newcommand{\sbbbr}[1]{\biggl[#1\biggr]}
\newcommand{\sbr}[1]{\left[#1\right]}
\newcommand{\half}{\frac{1}{2}}
\newcommand{\p}{\partial}
\newcommand{\jb}[1]{\langle#1\rangle}
\newcommand{\bset}[1]{\bigl\{#1\bigr\}}
\newcommand{\Bset}[1]{\Bigl\{#1\Bigr\}}
\newcommand{\set}[1]{\left\{#1\right\}}
\newcommand{\mytextfrac}[2]{{\textstyle\frac{#1}{#2}}}
\newcommand{\qaq}{\quad\text{and}\quad}
\newcommand{\para}{\parallel}
\newcommand{\ep}{\varepsilon}
\newcommand{\R}{\mathbb{R}}
\newcommand{\C}{\mathbb{C}}
\newcommand{\al}{\alpha}
\newcommand{\intRd}{\int_{\R^d}}
\newcommand{\intRdk}{\int_{\R^d_k}}
\newcommand{\intRdp}{\int_{\R^d_p}}
\newcommand{\lam}{\lambda}
\newcommand{\ceils}[1]{\lceil#1\rceil}
\newcommand{\Qin}{Q_{\mathrm{in}}}
\newcommand{\hin}{h_{\mathrm{in}}}
\newcommand\reallywidehat[1]{%
\savestack{\tmpbox}{\stretchto{%
\scaleto{%
\scalerel*[\widthof{\ensuremath{#1}}]{\kern.pt\mathchar"0362\kern.1pt}%
{\rule{0ex}{\textheight}}%WIDTH-LIMITED CIRCUMFLEX
}{\textheight}%
}{2.4ex}}%
\stackon[-6.9pt]{#1}{\tmpbox}%
}
\title{Small-parameter decay and uniform boundedness of the linearised operator symbol}
\author{Marnie Smith}
\date{doi:10.5802/jep.353; CC BY4.0}
\begin{document}
\maketitle
\noindent\textit{Editorial note.} The counted statement is the author's labelled lemma environment at source0.tex lines 1092-1106; its proof environment opens at line 1107, immediately after \textbackslash{}end\{lemma\}, so the pairing is by adjacency and the retained proof range 1107-1131 is the whole environment.\par\medskip
\noindent\textit{Editorial note.} Two environments are retained as context because the counted proof uses them: the directional integral estimate gnorm (source lines 1037-1042), used at source line 1108, and the labelled display Lingoodcoordinates (source lines 1059-1062), referred to at source lines 1108 and 1117. The three numbered items (a), (b), (c) used in the proof are declared inside the counted statement itself.\par\medskip
\noindent\textit{Editorial note.} This article carries no label-number dump and the retained counter values are left empty, so the printed numbering of the card is the wrapper's default; the identity of the counted item is its author label boundedness and decay lemma.\par\medskip
\noindent\textit{Editorial note.} The retained preamble is the publisher layout adapter, the native font and layout package and the native mathematical macro and theorem block of the pinned source.\par\medskip
\noindent\textit{Editorial scope.} One original result of the pinned publisher article, counted once: decay and uniform boundedness of the linearised operator symbol (author label boundedness and decay lemma).\par\medskip
\noindent\textit{Allowed standard prerequisites (editorial).}\begin{itemize}
\item Linearised operator symbol. The symbol of the linearised operator is given by an integral over the positive half-line of an exponential factor, a sine factor and the Fourier transform of the background.. It is the quantity whose decay and boundedness are asserted.
\item Sobolev norms with a weight. The norms appearing in the statement carry a derivative count, an extra positive power of one half and a logarithmic weight controlled by a parameter delta zero.. They are the norms that bound the symbol in the three parts of the statement.
\item Directional integral estimate. The directional integrals of derivatives of the background obey a bound by the weighted Sobolev norm, uniformly in the direction.. It is the estimate invoked in the first part of the proof.
\end{itemize}\medskip
\begin{lemma}\label{gnorm}
Let \( n \geq 0 \), and let \( \alpha \) be a multi-index. Then, for every $k\neq0$,
\begin{align*}
\int_0^\infty\jb{s}^{n}\babs{\p^\alpha_k\widehat g\bbr{\mytextfrac{k}{\abs{k}}s}}\,ds\lesssim_{d,
\delta_0,n,\aal}\norm{g}_{H^{n+{{1/2}}+\delta_0}_{\aal+\ceils{d/2}}}.\end{align*}
\end{lemma}

\begin{align}
\widetilde {\mathscr L}(\lam,k;\hbar)
= \frac{2}{\hbar\ak} \widehat w(k) \int_0^\infty e^{-\lam \sfrac{s}{\ak}} \sin(\mytextfrac{\hbar}{2} \ak s) \widehat g\bbr{\mytextfrac{k}{\ak}s} \, ds. \label{Lingoodcoordinates}
\end{align}

\begin{lemma}\label{boundedness and decay lemma}
Suppose $w\in L^1(\R^d)$ and $g\in H^{3/2+\delta_0}_{\ceils{d/2}}(\R^d)$. Then:\vspace*{-3pt}
\begin{enumeratea}
\item \label{(a)} As $\ak\to\infty$, the quantity $\widetilde {\mathscr L}(\lam,k;\hbar)\to0$, with decay uniform in $\Re\lam\geq0$ and $\hbar\in[0,1]$.

\item \label{(b)}There exists a constant $C_1=C_1(d,\delta_0)$, independent of $\hbar\in [0,1]$, such that for all $\Re\lam\geq0$ and $k\in\R^d$,\vspace*{-3pt}
\begin{align*}
\babs{\widetilde {\mathscr L}(\lam,k;\hbar)}\leq C_1\norm{w}_{L^1}\norm{g}_{H^{3/2+\delta_0}_{\ceils{d/2}}}.
\end{align*}
\item \label{(c)} If in addition $g\in H^{3/2+\delta_0}_{1+\ceils{d/2}}(\R^d)$, then there exists a constant $C_2=C_2(d,\delta_0)$, independent of $\hbar\in [0,1]$, such that for all $\Re\lam\geq0$ with $\lam\neq0$ and $k\in\R^d$,\vspace*{-3pt}
\begin{align*}
\babs{\widetilde {\mathscr L}(\lam,k;\hbar)}\leq C_2\frac{\ak}{\abs{\lam}}\norm{w}_{L^1}\norm{g}_{H^{3/2+\delta_0}_{1+\ceils{d/2}}}.
\end{align*}
\end{enumeratea}
\end{lemma}

\begin{proof}
First observe that \( \widetilde{\mathscr L}(\lam,k=0;\hbar) = 0 \) for all \( \Re \lam \geq 0 \) and \( \hbar \in [0,1] \). Now fix \( \hbar \in (0,1] \). For \( k \neq 0 \) and \( \Re \lam \geq 0 \), the representation \eqref{Lingoodcoordinates} together with Lemma~\ref{gnorm} yields the estimate\vspace*{-3pt}
\begin{align*}
\babs{\widetilde{\mathscr L}(\lam,k;\hbar)}
&\lesssim_d \abs{\widehat w(k)} \frac{1}{\hbar \ak} \int_0^\infty \babs{\sin\br{\mytextfrac{\hbar}{2} \ak s} \, \widehat g\bbr{\mytextfrac{k}{\ak}s}}\,ds \\
&\lesssim \abs{\widehat w(k)} \int_0^\infty s\, \babs{\widehat g\bbr{\mytextfrac{k}{\ak}s}}\,ds \\
&\lesssim_{d,\delta_0} \abs{\widehat w(k)} \norm{g}_{H^{3/2+\delta_0}_{\ceils{d/2}}}.
\end{align*}
Since \( w \in L^1(\R^d) \), the Fourier transform \( \widehat w(k) \to 0 \) as \( \abs{k} \to \infty \), establishing part~\eqref{(a)}. Moreover, \( \abs{\widehat w(k)} \leq \norm{w}_{L^1} \) for all \( k \), yielding part~\eqref{(b)}.

For part~\eqref{(c)}, the expression \eqref{Lingoodcoordinates} can be integrated by parts in \( s \), with no boundary contributions, to obtain\vspace*{-1pt}
\begin{align*}
\widetilde{\mathscr L}&(\lam,k;\hbar)
\\[-3pt]
&= 2 \widehat w(k) \frac{1}{\hbar \ak} \Bigl(-\frac{\ak}{\lam}\Bigr) \int_0^\infty \p_s \bigl( e^{-\lam \sfrac{s}{\ak}} \bigr) \sin(\mytextfrac{\hbar}{2} \ak s) \, \widehat g\bbr{\mytextfrac{k}{\ak}s} \, ds \\
&= 2 \widehat w(k) \frac{1}{\hbar \lam} \int_0^\infty e^{-\lam \sfrac{s}{\ak}} \left[ \mytextfrac{\hbar}{2} \ak \cos(\mytextfrac{\hbar}{2} \ak s) \widehat g\bbr{\mytextfrac{k}{\ak}s} + \sin(\mytextfrac{\hbar}{2} \ak s) \mytextfrac{k}{\ak} \cdot \nabla_k \widehat g\bbr{\mytextfrac{k}{\ak}s} \right] ds.
\end{align*}
Therefore,\vspace*{-5pt}
\begin{align*}
\abs{\widetilde{\mathscr L}(\lam,k;\hbar)}
&\lesssim \frac{\ak}{\abs{\lam}} \norm{w}_{L^1} \int_0^\infty \Bigl[ \babs{\widehat g\bbr{\mytextfrac{k}{\ak}s}} + s \babs{\nabla_k \widehat g\bbr{\mytextfrac{k}{\ak}s}} \Bigr] ds \\
&\lesssim_{d,\delta_0} \frac{\ak}{\abs{\lam}} \norm{w}_{L^1} \norm{g}_{H^{3/2+\delta_0}_{1+\ceils{d/2}}}.
\end{align*}
The same bounds hold for the classical limit \( \widetilde{\mathscr L}(\lam,k;0) \), with identical estimates.
\end{proof}

\begin{thebibliography}{99}
\bibitem{r1} J. Bedrossian - “Nonlinear echoes and Landau damping with insufficient regularity”, Tunis. J. Math. 3 (2021) no. 1, p. 121-205
\bibitem{r2} J. Bedrossian, N. Masmoudi \& C. Mouhot - “Landau damping: paraproducts and Gevrey regularity”, Ann. PDE 2 (2016) no. 1, article ID 4, 71 pages
\bibitem{r3} J. Bedrossian, N. Masmoudi \& C. Mouhot - “Landau damping in finite regularity for unconfined systems with screened interactions”, Comm. Pure Appl. Math. 71 (2018) no. 3, p. 537-576
\bibitem{r4} J. Bedrossian, N. Masmoudi \& C. Mouhot - “Linearized wave-damping structure of Vlasov–Poisson in $\mathbb{R}^3$”, SIAM J. Math. Anal. 54 (2022) no. 4, p. 4379-4406
\bibitem{r5} A. Borie, S. Hadama \& J. Sabin - “Scattering for the positive density Hartree equation”, 2025
\bibitem{r6} A. Bove, G. Da Prato \& G. Fano - “An existence proof for the Hartree-Fock time-dependent problem with bounded two-body interaction”, Comm. Math. Phys. 37 (1974), p. 183-191
\bibitem{r7} A. Bove, G. Da Prato \& G. Fano - “On the Hartree-Fock time-dependent problem”, Comm. Math. Phys. 49 (1976) no. 1, p. 25-33
\bibitem{r8} T. Cazenave - Semilinear Schrödinger equations, Courant Lecture Notes in Math., vol. 10, American Mathematical Society, Providence, RI, 2003
\bibitem{r9} J. M. Chadam - “The time-dependent Hartree-Fock equations with Coulomb two-body interaction”, Comm. Math. Phys. 46 (1976) no. 2, p. 99-104
\bibitem{r10} T. Chen, Y. Hong \& N. Pavlović - “Global well-posedness of the NLS system for infinitely many fermions”, Arch. Rational Mech. Anal. 224 (2017) no. 1, p. 91-123
\bibitem{r11} T. Chen, Y. Hong \& N. Pavlović - “On the scattering problem for infinitely many fermions in dimensions $d\ge 3$ at positive temperature”, Ann. Inst. H. Poincaré C Anal. Non Linéaire 35 (2018) no. 2, p. 393-416
\bibitem{r12} J. J. Chong, L. Lafleche \& C. Saffirio - “On the $L^2$ rate of convergence in the limit from the Hartree to the Vlasov-Poisson equation”, J. Éc. polytech. Math. 10 (2023), p. 703-726
\bibitem{r13} C. Collot \& A.-S. de Suzzoni - “Stability of equilibria for a Hartree equation for random fields”, J. Math. Pures Appl. (9) 137 (2020), p. 70-100
\bibitem{r14} C. Collot \& A.-S. de Suzzoni - “Stability of steady states for Hartree and Schrödinger equations for infinitely many particles”, Ann. H. Lebesgue 5 (2022), p. 429-490
\bibitem{r15} E. Grenier, T. T. Nguyen \& I. Rodnianski - “Landau damping for analytic and Gevrey data”, Math. Res. Lett. 28 (2021) no. 6, p. 1679-1702
\bibitem{r16} S. Hadama - “Asymptotic stability of a wide class of stationary solutions for the Hartree and Schrödinger equations for infinitely many particles”, Ann. H. Lebesgue 8 (2025), p. 181-218
\bibitem{r17} A. D. Ionescu, B. Pausader, X. Wang \& K. Widmayer - “Nonlinear Landau damping for the Vlasov-Poisson system in $\mathbb{R}^3$: the Poisson equilibrium”, Ann. PDE 10 (2024) no. 1, article ID 2, 78 pages
\bibitem{r18} L. Landau - “On the vibrations of the electronic plasma”, Acad. Sci. USSR. J. Phys. 10 (1946), p. 25-34
\bibitem{r19} M. Lewin \& J. Sabin - “The Hartree equation for infinitely many particles, II: Dispersion and scattering in 2D”, Anal. PDE 7 (2014) no. 6, p. 1339-1363
\bibitem{r20} M. Lewin \& J. Sabin - “The Hartree equation for infinitely many particles I. Well-posedness theory”, Comm. Math. Phys. 334 (2015) no. 1, p. 117-170
\bibitem{r21} M. Lewin \& J. Sabin - “The Hartree and Vlasov equations at positive density”, Comm. Partial Differential Equations 45 (2020) no. 12, p. 1702-1754
\bibitem{r22} C. Mouhot \& C. Villani - “On Landau damping”, Acta Math. 207 (2011) no. 1, p. 29-201
\bibitem{r23} T. T. Nguyen \& C. You - “Modified scattering for long-range Hartree equations of infinite rank near vacuum”, SIAM J. Math. Anal. 57 (2025) no. 6, p. 6486-6497
\bibitem{r24} T. T. Nguyen \& C. You - “Plasmons for the Hartree equations with Coulomb interaction”, Probab. Math. Phys. 6 (2025) no. 3, p. 913-960
\bibitem{r25} E. M. Stein \& G. Weiss - “Interpolation of operators with change of measures”, Trans. Amer. Math. Soc. 87 (1958), p. 159-172
\bibitem{r26} C. You - “Phase mixing estimates for the nonlinear Hartree equation of infinite rank”, 2024
\end{thebibliography}
\end{document}
